\section{How the constants are computed and certified}
\label{sec:computation}

\subsection{An exact finite combinatorial calculation}

The common divisor does not require numerical integration over a
circle. The calculation is finite:
\begin{enumerate}
\item Form the integer slopes \(0,t_j,L-t_j\) and all nonzero
differences.
\item Form and sort the rational critical points \(a/d\in[0,1]\).
\item In each open critical cell choose its rational midpoint \(z\).
Sort all fractional endpoints \(\{t_jz\},\{(L-t_j)z\}\).
\item Evaluate the integer vector \(W(z,y)\) at each rational
midpoint between consecutive endpoints. Minimize its weighted sum.
\item Apply the finite digamma sum~\eqref{eq:gamma}. Adjacent cells
with the same minimum may be merged.
\end{enumerate}
Only the final special-function evaluations use real arithmetic.
All endpoint orderings and minima use exact fractions and integers.

For the base weights \((6,1,\ldots,1)\) at \(L=69\), the first
merged nonzero cells are:
\begin{center}
\begin{tabular}{ccc}
\toprule
Left endpoint&Right endpoint&\(f\)\\
\midrule
\(1/69\)&\(1/67\)&6\\
\(1/67\)&\(1/65\)&7\\
\(1/65\)&\(1/62\)&8\\
\(1/62\)&\(1/60\)&7\\
\(1/60\)&\(1/58\)&6\\
\(1/58\)&\(1/56\)&5\\
\(1/56\)&\(1/54\)&4\\
\(1/54\)&\(1/52\)&3\\
\bottomrule
\end{tabular}
\end{center}
These are only the beginning of the partition. Endpoint wraparound
and changing overlaps make a full computation preferable to
extrapolating a pattern from the first few cells.
The complete merged list is supplied in the certificate data.

\subsection{Ball arithmetic, root signs, and strict rounding}

The standalone program \texttt{code/certify.py} requires Python and
\texttt{python-flint}. It uses Arb~\cite{Arb} at 256-bit precision.
A ball encloses the exact real value by a midpoint and a rigorously
propagated error radius.

The saddle equation is strictly decreasing. At \(Y=1/2\) its
left side minus \(\Theta\) is positive, and at \(Y=2M\) it is
negative. The program performs 210 rational bisections, accepting
each sign only if the entire corresponding ball has that sign.
It then evaluates \(D_0,G_0,\gamma,D,G,\Delta_k,E_k\) on the
resulting interval and asserts both \(\Delta_k>0\) and the advertised
strict upper bound.

For example, in the twelve-layer cubic case it encloses
\[
E_3=1203.97243920634698117089090450914656\ldots
\]
with an error radius below \(6\cdot10^{-59}\).
The assertion \(E_3<1203.973\) has a wide safety margin relative
to this radius. The certificate does not assume that a floating-point
optimizer located a true minimum.

The same program checks the arc bound independently against
all integer pole positions for \(n=345\): 52 nonexceptional primes
are checked, and 39 have a positive guaranteed divisor exponent.
This finite check tests the combinatorial implementation; the
proof for arbitrary \(n\) is Lemma~\ref{lem:valuation}.

\subsection{Reproduction and the role of each file}

Run the following from the project directory:
\begin{verbatim}
python code/certify.py
python code/finite_checks.py
xelatex -interaction=nonstopmode -halt-on-error main.tex
xelatex -interaction=nonstopmode -halt-on-error main.tex
\end{verbatim}
The first command regenerates \texttt{data/certificates.json}.
The second requires SymPy and checks the small exact examples in
Appendix~\ref{app:examples}.
Optional experiment programs require the additional packages stated
in the README. They are not prerequisites for the main proof.

The certificate records the parameters, merged nonzero arc cells,
root enclosure, all rates, the gap, and the strict comparisons.
Its purpose is to make the numerical part reproducible.
The prime number theorem and the asymptotic inequalities are
mathematical inputs, not claims checked by a finite program.

\subsection{Parameter order and effectivity}

For a chosen odd \(k\), first fix \(L,t_j,\nu_j\), or fix the scale
\(s\) in the twelve-layer family. Then fix an exponent slack
\(\varepsilon\), and choose the small \(\eta\) in
Section~\ref{sec:transfer}. Only after these choices does \(n\)
tend to infinity, and only then is \(n\) selected as a function of
the approximation denominator \(q\).

This manuscript gives explicit upper bounds on the irrationality
\emph{exponent}. It does not tabulate an explicit numerical
\(q_0(\varepsilon)\) or the constants in~\eqref{eq:measuremeaning}.
Those would require explicit versions of the prime-counting
remainder and all finite error terms. They are a separate
effectivity calculation, not part of the decimal certification.
